\clearpage
\section{Второй промежуточный отчёт}

Второй этап учебной практики был посвящён механизмам, позволяющим программе исследовать собственную структуру, загружать внешние компоненты и формировать программный код во время выполнения. В течение недели были последовательно выполнены пять заданий: «Атрибуты классов и рефлексия в C\#», «Динамические библиотеки», «Метаданные библиотеки классов», «Система плагинов» и «Генерация классов во время выполнения программы». Все решения включены в общее решение \texttt{practice2025}, размещены в отдельных ветках репозитория и проверены средствами непрерывной интеграции.

Выполненные работы образуют логически связанную последовательность. Сначала были исследованы пользовательские атрибуты как способ декларативного описания программных сущностей. Затем полученные знания применялись при динамической загрузке сборок и извлечении их метаданных. На следующем этапе эти механизмы были объединены в систему плагинов с учётом зависимостей. Завершающее задание показало обратную сторону динамического программирования: программа не только анализировала готовые типы, но и создавала новый тип из исходного текста. Согласно документации Microsoft, атрибуты добавляют к программным сущностям метаданные, доступные во время выполнения через рефлексию \cite{attributes}. Это положение стало общей теоретической основой большинства заданий второй недели.

\begin{table}[h]
    \centering
    \caption{Состав выполненных работ второй недели}
    \small
    \begin{tabular}{>{\centering\arraybackslash}p{0.8cm} >{\RaggedRight\arraybackslash}p{8.5cm} >{\centering\arraybackslash}p{2cm} >{\centering\arraybackslash}p{2.2cm}}
        \toprule
        \textbf{№} & \textbf{Задание} & \textbf{Тестов} & \textbf{Результат CI} \\
        \midrule
        7 & Атрибуты классов и рефлексия в C\# & 6 & успешно \\
        8 & Динамические библиотеки & 4 & успешно \\
        9 & Метаданные библиотеки классов & 4 & успешно \\
        10 & Система плагинов & 6 & успешно \\
        11 & Генерация классов во время выполнения & 8 & успешно \\
        \bottomrule
    \end{tabular}
\end{table}

\subsection{Атрибуты классов и рефлексия в C\#}

Выполнение задания началось с определения назначения пользовательских атрибутов. Было установлено, что атрибут представляет собой декларативную характеристику типа или его члена и не изменяет поведение программы самостоятельно. Реакция на атрибут появляется только в том компоненте, который считывает метаданные и интерпретирует их. На основании этого положения были созданы классы \texttt{DisplayNameAttribute} и \texttt{VersionAttribute}. Для каждого атрибута была задана допустимая область применения. Атрибут отображаемого имени применялся к классам, методам и свойствам, а атрибут версии --- к классу.

После определения атрибутов был создан демонстрационный класс \texttt{SampleClass}. Класс получил отображаемое имя «Пример класса» и версию 1.0. Его метод и свойство также были отмечены атрибутами отображаемого имени. Затем был реализован статический класс \texttt{ReflectionHelper}. Метод \texttt{PrintTypeInfo} принимал объект \texttt{Type}, считывал атрибуты класса и перебирал его публичные методы и свойства. Результат выводился в консоль в понятном человеку виде. Корректность решения проверялась шестью модульными тестами: четыре теста подтверждали значения отдельных атрибутов, ещё два проверяли полный вывод и поведение для типа без пользовательских атрибутов.

Основная трудность состояла в разграничении стандартных и пользовательских членов типа. При получении методов через рефлексию в результат могут попадать унаследованные методы \texttt{Object}, а свойства имеют методы доступа, которые также присутствуют в метаданных. Проблема была решена явным выбором требуемых элементов и проверкой наличия \texttt{DisplayNameAttribute} перед формированием строки вывода. Такой подход позволил не включать в результат служебную информацию, не относящуюся к заданию.

Наиболее сложным моментом стало понимание связи между декларативной записью атрибута и выполняемой логикой. Внешне атрибут напоминает специальную команду, однако фактически он является объектом метаданных. Поэтому требовалось отдельно реализовать код, который обнаруживает атрибут и использует его свойства. Данная особенность была закреплена тестом для класса без атрибутов: отсутствие метаданных не должно приводить к ошибке и не должно создавать вымышленные значения.

\subsection{Динамические библиотеки}

В задании о динамических библиотеках были созданы три функциональные части. Первая часть представлена библиотекой \texttt{CommandLib}, содержащей интерфейс \texttt{ICommand} с методом \texttt{Execute}. Вторая часть --- библиотека \texttt{FileSystemCommands} с двумя реализациями команд. Команда \texttt{DirectorySizeCommand} вычисляет суммарный размер файлов выбранного каталога, а \texttt{FindFilesCommand} находит файлы, соответствующие переданной маске. Третьей частью стало консольное приложение \texttt{CommandRunner}, которое получает путь к библиотеке, загружает сборку, обнаруживает реализации интерфейса, создаёт их экземпляры и выполняет команды.

Для загрузки использовался метод \texttt{Assembly.LoadFrom}, предназначенный для загрузки сборки по указанному пути \cite{loadfrom}. После загрузки выполнялся анализ типов сборки. Кандидат должен был быть конкретным классом, реализующим \texttt{ICommand}. Создание экземпляра выполнялось с учётом параметров конструктора конкретной команды. Такой подход продемонстрировал, что интерфейс служит стабильным контрактом между приложением и динамически подключаемой библиотекой.

Трудности возникли при тестировании команд, выводящих результат в консоль. Первоначальная проверка только отсутствия исключения не доказывала правильность вычисления. Для повышения точности тестов стандартный поток вывода временно перенаправлялся в \texttt{StringWriter}. После выполнения команды из полученной строки извлекался результат и сравнивался с ожидаемым значением. Для каждого теста создавался отдельный временный каталог с заранее известными файлами, а после проверки каталог удалялся. Это обеспечило независимость тестов от содержимого пользовательской файловой системы.

Наиболее сложной частью стала организация границы между основной программой и загружаемой библиотекой. Если интерфейс объявить только в приложении или продублировать в разных проектах, одинаковые по имени интерфейсы будут различными типами для среды выполнения. Решением стало вынесение контракта в отдельную библиотеку \texttt{CommandLib}, на которую ссылаются обе стороны. В результате консольное приложение может работать с неизвестными на этапе компиляции реализациями, если они соблюдают общий контракт.

\subsection{Метаданные библиотеки классов}

Задание по метаданным продолжило работу с динамическими сборками, но изменило цель анализа. Вместо поиска только команд требовалось вывести полное структурированное описание библиотеки: список классов, их атрибутов, методов и конструкторов, а также имена и типы параметров методов и конструкторов. Для решения было разработано консольное приложение, принимающее путь к DLL через аргументы командной строки.

После проверки существования файла сборка загружалась в память. Для каждого класса последовательно формировались разделы с атрибутами, конструкторами и методами. Параметры преобразовывались в текстовые записи, включающие имя типа и имя параметра. Для методов дополнительно выводился тип возвращаемого значения. Атрибуты из задания 07 были добавлены к анализируемым классам, благодаря чему программа демонстрировала не только стандартные метаданные .NET, но и пользовательские значения.

Основное затруднение было связано с объёмом данных, доступных через рефлексию. В сборке присутствуют автоматически созданные методы доступа к свойствам, унаследованные члены и служебные типы. Неконтролируемый вывод делал результат громоздким и затруднял его проверку. Проблема была решена с помощью явных условий отбора и единообразных методов форматирования. Для сохранения предсказуемого порядка элементы сортировались по имени. Это особенно важно для модульных тестов, поскольку один и тот же набор метаданных должен формировать стабильный текст на разных запусках.

Наиболее сложным моментом оказалось корректное представление сигнатур. Метод и конструктор могут не иметь параметров, содержать несколько параметров различных типов или использовать типы с пространствами имён. Было необходимо избежать неоднозначных строк и одновременно сохранить читаемость результата. Решение состояло в общем преобразовании каждого \texttt{ParameterInfo} в пару «тип --- имя» и объединении таких пар в порядке следования параметров. Корректность была подтверждена тестами для классов, атрибутов, методов и конструкторов.

\subsection{Система плагинов}

В задании о системе плагинов были объединены интерфейсы, атрибуты, динамическая загрузка и анализ зависимостей. В качестве общего контракта использовался интерфейс \texttt{ICommand} из задания 08. Был определён атрибут \texttt{PluginLoadAttribute}, которым отмечались классы, предназначенные для автоматической загрузки. Второй атрибут, \texttt{PluginDependencyAttribute}, хранил тип плагина, который должен быть загружен раньше текущего.

Реализованный \texttt{PluginManager} обнаруживает DLL-файлы в указанном каталоге, загружает сборки и выбирает конкретные классы, одновременно реализующие \texttt{ICommand}, имеющие открытый конструктор без параметров и отмеченные атрибутом \texttt{PluginLoad}. Затем строится порядок загрузки. В демонстрационном наборе плагин регистратора не имеет зависимостей, хранилище зависит от регистратора, а генератор отчётов зависит от хранилища. После определения порядка создаются экземпляры и последовательно вызывается метод \texttt{Execute}.

Главная трудность была связана с тем, что простой порядок файлов в каталоге не отражает логических зависимостей. Если загружать плагины в алфавитном или случайном порядке, зависимый компонент может начать работу раньше необходимого ему компонента. Зависимости были представлены ориентированным графом. Для каждого плагина выполнялся поиск в глубину: сначала обрабатывались зависимости, после чего в результат добавлялся сам плагин. Дополнительно использовались состояния «обрабатывается» и «обработан». Повторное посещение вершины в состоянии «обрабатывается» означает цикл, при котором корректный порядок не существует.

Наиболее сложным моментом стала обработка ошибочных конфигураций. Отсутствующую зависимость нельзя игнорировать, поскольку тогда система запустит плагин в недопустимом состоянии. Циклическую зависимость также нельзя разрешить выбором произвольного порядка. Поэтому для обоих случаев формируются явные исключения с понятным описанием причины. Шесть тестов проверяют обнаружение всех плагинов, наличие атрибутов, создание экземпляров, правильный порядок выполнения и выявление цикла.

\subsection{Генерация классов во время выполнения программы}

Заключительное практическое задание состояло в создании класса \texttt{Calculator}, отсутствующего в скомпилированной библиотеке. Текст класса хранился в строковой константе и включал четыре метода: \texttt{Add}, \texttt{Minus}, \texttt{Mul} и \texttt{Div}. Для компиляции использовались API платформы Roslyn, предоставляющие программный доступ к синтаксической и семантической модели C\# \cite{roslyn}.

Сначала исходная строка преобразовывалась в синтаксическое дерево. Затем формировался набор ссылок на системные сборки и сборку с контрактом \texttt{ICalculator}. Объект \texttt{CSharpCompilation} создавал динамическую библиотеку в потоке памяти. При неуспешной компиляции диагностические сообщения объединялись и включались в исключение. При успешной компиляции байты сборки загружались в память, находился тип \texttt{Calculator} и создавался его экземпляр.

Условие запрещало вызывать арифметические методы через рефлексию. Проблема была решена изменением исходного определения класса: динамический \texttt{Calculator} реализует заранее известный интерфейс \texttt{ICalculator}. Рефлексия применяется только на границе создания экземпляра, поскольку имя нового типа неизвестно скомпилированному коду. После приведения объекта к интерфейсу выражения \texttt{calculator.Add(a, b)} и аналогичные вызовы являются обычными виртуальными вызовами через контракт и не используют \texttt{MethodInfo.Invoke}.

Трудность заключалась в формировании полного набора ссылок компиляции. Наличие типа \texttt{object} не гарантирует доступность всех системных типов, используемых компилятором. Для решения использовался список доверенных платформенных сборок текущей среды выполнения, к которому добавлялась сборка с интерфейсом \texttt{ICalculator}. Наиболее сложным оказался переход между двумя мирами типов: статически известным интерфейсом и типом, появившимся только после компиляции строки. Восемь тестовых методов проверили арифметические операции, деление на ноль, создание экземпляра, использование переданной строки и обработку синтаксической ошибки.

\subsection{Задачи генерации во время компиляции и выполнения}

Генерация кода во время компиляции решает задачи устранения повторяющегося шаблонного кода, построения адаптеров, сериализаторов, регистрационных таблиц и иных структур, которые могут быть полностью определены по исходной программе. Результат такой генерации становится частью итоговой сборки. Компилятор проверяет типы, а разработчик получает ошибки до запуска приложения. Подход целесообразен, когда входные данные известны заранее и редко изменяются во время работы системы. Примерами являются генерация клиентов по описанию API, реализация преобразований данных и автоматическое создание вспомогательных методов для моделей.

Генерация во время выполнения применяется тогда, когда исходная информация становится известна только после запуска. К ней относятся пользовательские формулы, сценарии, динамические правила обработки и специализированные запросы. Программа может получить текст или модель операции, скомпилировать её и выполнить без выпуска новой версии основного приложения. Такой подход повышает гибкость, но переносит часть ошибок с этапа сборки на этап выполнения и требует дополнительных мер контроля.

Различие между подходами определяется моментом фиксации результата. Компиляционная генерация увеличивает объём работы сборочной системы, зато созданный код может быть заранее проанализирован и протестирован. Генерация во время выполнения допускает адаптацию к текущим данным и настройкам, но потребляет время и память при запуске, усложняет диагностику и может создавать риск исполнения непроверенного кода. Следовательно, выбор должен основываться не на максимальной динамичности, а на том, действительно ли изменяемость после поставки является требованием системы.

\subsection{Статическая и динамическая компоновка}

При статической компоновке состав зависимостей определяется во время сборки. Проект содержит ссылки на библиотеки, компилятор проверяет доступность типов и методов, а необходимые компоненты поставляются как заранее определённый набор. Изменение контракта обычно выявляется ошибкой компиляции. Такой подход обеспечивает высокую предсказуемость, упрощает навигацию по коду и позволяет инструментам разработки выполнять полный статический анализ.

Динамическая компоновка переносит выбор конкретной реализации на время выполнения. Основное приложение может знать только общий интерфейс, а реализация обнаруживается по пути к DLL, конфигурации или атрибуту. Это позволяет добавлять компоненты без повторной компиляции ядра, выбирать разные реализации для разных условий и уменьшать прямую связанность модулей. На этой основе строятся расширяемые редакторы, серверные платформы и системы обработки данных.

Преимуществами статической компоновки являются раннее обнаружение ошибок, простая модель развёртывания и меньшая неопределённость поведения. Недостатком является необходимость пересборки и повторной поставки приложения при изменении набора компонентов. Динамическая компоновка обеспечивает расширяемость и независимое обновление модулей, но создаёт риск несовместимости версий, отсутствия файла, ошибки загрузки и нарушения порядка инициализации. Кроме того, часть связей становится невидимой для компилятора и должна проверяться собственным кодом и интеграционными тестами.

Рациональная архитектура часто сочетает оба подхода. Базовые функции и контракты связываются статически, поскольку они определяют устойчивое ядро приложения. Дополнительные возможности подключаются динамически через узкие интерфейсы. В выполненной работе статической частью являлся \texttt{ICommand}, а плагины регистратора, хранилища и генератора отчётов представляли динамические реализации.

\subsection{Влияние системы плагинов на архитектуру}

Система плагинов переводит архитектуру от монолитного набора функций к модели «ядро --- расширения». Ядро отвечает за жизненный цикл, обнаружение, проверку и запуск компонентов. Плагин реализует ограниченный контракт и не должен зависеть от внутренних деталей приложения. Такое разделение уменьшает необходимость изменять ядро при добавлении новых возможностей и поддерживает принцип открытости для расширения.

Одновременно архитектура объективно усложняется. Появляются правила совместимости, каталогизация расширений, обработка ошибок загрузки, управление версиями и порядок инициализации. При наличии зависимостей требуется граф, а не простой список. Становится сложнее выполнять сквозную отладку, поскольку часть кода может быть поставлена независимо от основной программы. Поэтому сама по себе система плагинов не является упрощением; она обменивает простоту небольшой программы на управляемую расширяемость большой системы.

Усложнение оправдано, если приложение действительно должно поддерживать независимые команды, интеграции или функциональные модули. Для небольшой программы с фиксированным набором операций плагины создадут избыточную инфраструктуру. Для редактора, среды разработки или корпоративной платформы преимущества независимого развития компонентов обычно превышают расходы на механизм загрузки. Важным условием является минимальный и стабильный контракт: чем больше внутреннего состояния ядра доступно плагину, тем сильнее связанность и тем труднее обновление.

\subsection{Риски безопасности и способы их снижения}

Основной риск плагинной системы связан с тем, что динамически загруженная библиотека является исполняемым кодом. Получив те же права, что и основной процесс, вредоносный плагин может читать и изменять файлы пользователя, отправлять данные по сети, запускать процессы или нарушать работу приложения. Цифровая подпись файла сама по себе не делает код безопасным, однако позволяет проверить издателя и целостность поставки.

Второй риск состоит в подмене библиотеки. Если каталог плагинов доступен для записи постороннему процессу или путь задаётся без проверки, приложение может загрузить другой файл с подходящим именем. Третий риск связан с уязвимостями обычного добросовестного плагина: ошибка обработки входных данных, чрезмерное потребление памяти или бесконечное выполнение способны привести к отказу всей программы. Дополнительную опасность создают несовместимые зависимости и загрузка нескольких версий одной библиотеки.

Снижение рисков требует сочетания организационных и технических мер. Следует загружать расширения только из доверенного каталога, проверять издателя и контрольную сумму, ограничивать перечень разрешённых плагинов и журналировать операции загрузки. Контракт должен предоставлять минимально необходимые возможности. Потенциально опасные расширения целесообразно запускать в отдельном процессе с ограниченными правами, поскольку обычная загрузка DLL внутрь процесса не создаёт границы безопасности. Для внешнего процесса можно установить ограничения времени, памяти и доступа к ресурсам, а взаимодействие организовать через проверяемый протокол.

Необходимо также обрабатывать исключения на границе плагина, проверять совместимость версий и не передавать расширению секретные данные без необходимости. При динамической генерации кода аналогичные риски возникают при компиляции текста, поступившего от пользователя. Запуск такого кода в основном процессе допустим только для полностью доверенного источника. В противном случае требуется изоляция или отказ от произвольного языка в пользу ограниченной предметно-ориентированной модели.

\subsection{Примеры систем плагинов}

С системами плагинов пользователь регулярно сталкивается в веб-браузерах. Расширения могут изменять страницы, блокировать рекламу, управлять паролями и добавлять средства разработчика. Браузер предоставляет им формализованный API и запрашивает разрешения на доступ к вкладкам, сайтам и данным. Этот пример показывает одновременно преимущества расширяемости и необходимость модели безопасности.

Другим примером являются среды разработки Visual Studio и Visual Studio Code. Расширения добавляют поддержку языков, форматирование, анализ кода, интеграцию с системами контроля версий и пользовательские команды. Ядро среды не обязано содержать реализацию каждого языка, поскольку функциональность развивается независимыми поставщиками. Аналогичная модель используется в графических редакторах и программах трёхмерного моделирования, где плагины реализуют форматы файлов, фильтры и инструменты обработки.

Системы управления содержимым также применяют плагины. В WordPress расширения добавляют формы, электронную коммерцию, оптимизацию и интеграции. Игровые проекты используют модификации для новых объектов, сценариев и правил. Медиапроигрыватели подключают кодеки и модули вывода. Во всех случаях сохраняется общий принцип: ядро определяет точки расширения, а внешние компоненты реализуют договорённое поведение.

Практические задания показали внутреннее устройство подобных решений в уменьшенном масштабе. Атрибут \texttt{PluginLoad} выполнял роль декларации расширения, интерфейс \texttt{ICommand} задавал точку взаимодействия, рефлексия обеспечивала обнаружение, а граф зависимостей определял порядок запуска. Поэтому знакомые пользовательские системы перестали восприниматься как единый неразделимый продукт: их функциональность может быть результатом совместной работы ядра и множества независимо загружаемых модулей.

\subsection{Итоги рефлексивного анализа второй недели}

Работы второй недели продемонстрировали последовательный переход от статически определённой структуры программы к управляемой динамичности. Атрибуты позволили описывать намерение, рефлексия --- считывать это описание, динамическая загрузка --- подключать неизвестные заранее сборки, система плагинов --- организовывать их совместную работу, а Roslyn --- создавать новый тип из исходного текста.

Наиболее значимые трудности были связаны с границами между компонентами. Для динамических библиотек такой границей являлся общий интерфейс, для метаданных --- правила отбора и форматирования, для плагинов --- граф зависимостей, а для динамической генерации --- интерфейс между статически скомпилированным кодом и новым типом. Во всех случаях устойчивое решение достигалось за счёт явного контракта, проверки ошибочных состояний и модульных тестов.

Полученные результаты подтверждают, что динамические механизмы не отменяют требования строгой структуры. Напротив, чем позже определяется конкретная реализация, тем точнее должны быть заданы контракты, правила совместимости и меры безопасности. Следовательно, профессиональное применение рефлексии, плагинов и генерации кода требует не только знания API, но и архитектурного анализа последствий каждого динамического решения.
